\documentclass[12pt]{article}
\usepackage[utf8]{inputenc}
\usepackage[russian]{babel}
\usepackage{amsthm,amsmath,amssymb}
\usepackage{fullpage}

\newtheorem{theorem}{Теорема}
\newtheorem{lemma}{Лемма}

\begin{document}

  \title{Продвинутые алгоритмы: алгоритмы для NP‑трудных задач}
  \author{Кожагулов Рамазан}
  \date{}
  \maketitle

  \section*{Задача 1.1. Покрытие гласных с минимальной суммарной длиной}

  \paragraph{Условие.}
  Из слов \emph{сыроежка}, \emph{опёнок}, \emph{трюфель}, \emph{шампиньон}, \emph{мухомор}, \emph{белянка}, \emph{рыжик}, \emph{эвкалипт}, \emph{берёза}, \emph{туя}, \emph{секвойя}, \emph{липа}, \emph{крыжовник} требуется выбрать набор так, чтобы среди всех выбранных слов встретились все русские гласные буквы \\ \(U=\{\text{а},\text{е},\text{ё},\text{и},\text{о},\text{у},\text{ы},\text{э},\text{ю},\text{я}\}\), а суммарное число букв в выбранных словах было минимальным. Сформулировать задачу как задачу о покрытии и решить Алгоритмом
Хватала

  \paragraph{Модель как задача о покрытии множествами (Set Cover).}
  Для каждого слова \(w\) из списка определим подмножество \(S_w\subseteq U\), состоящее из различных гласных, встречающихся в \(w\). Вес множества равен длине слова: \(c(w)=|w|\). Требуется найти множество слов \(\mathrm{Sol}\) минимальной суммарной стоимости \(\sum_{w\in\mathrm{Sol}} c(w)\) такое, что \(\bigcup_{w\in\mathrm{Sol}} S_w = U\). 

  Для дальнейшего удобно выписать \(S_w\) и \(c(w)\) для всех слов:
  \begin{align*}
    c(\text{сыроежка}) &= 8, \quad S = \{\text{а},\text{е},\text{о},\text{ы}\}; \\
    c(\text{опёнок}) &= 6, \quad S = \{\text{о},\text{ё}\}; \\
    c(\text{трюфель}) &= 7, \quad S = \{\text{ю},\text{е}\}; \\
    c(\text{шампиньон}) &= 9, \quad S = \{\text{а},\text{и},\text{о}\}; \\
    c(\text{мухомор}) &= 7, \quad S = \{\text{у},\text{о}\}; \\
    c(\text{белянка}) &= 7, \quad S = \{\text{е},\text{я},\text{а}\}; \\
    c(\text{рыжик}) &= 5, \quad S = \{\text{ы},\text{и}\}; \\
    c(\text{эвкалипт}) &= 8, \quad S = \{\text{э},\text{а},\text{и}\}; \\
    c(\text{берёза}) &= 6, \quad S = \{\text{е},\text{ё},\text{а}\}; \\
    c(\text{туя}) &= 3, \quad S = \{\text{у},\text{я}\}; \\
    c(\text{секвойя}) &= 7, \quad S = \{\text{е},\text{о},\text{я}\}; \\
    c(\text{липа}) &= 4, \quad S = \{\text{и},\text{а}\}; \\
    c(\text{крыжовник}) &= 9, \quad S = \{\text{ы},\text{о},\text{и}\}.
  \end{align*}

  \paragraph{Применение алгоритма Хватала.}
  Будем итеративно выбирать слово \(w\), минимизирующее \(\alpha(w)=\dfrac{c(w)}{|S_w\setminus C|}\), где \(C\) — множество уже покрытых гласных. При равенстве значений \(\alpha\) выбираем слово, покрывающее больше новых гласных, а при дальнейшем равенстве — первое по порядку в приведённом списке.

  \textbf{1-я итерация.} \(C=\varnothing\). Минимум достигается у слова \emph{туя} с \(\alpha=3/2=1{,}5\). Добавляем \emph{туя}, получаем \(C=\{\text{у},\text{я}\}\), стоимость \(3\).

  \textbf{2-я итерация.} Рассматриваем только непокрытые гласные. Для \emph{сыроежка} имеем \(\alpha=8/4=2\), для \emph{берёза} \(\alpha=6/3=2\), для \emph{липа} \(\alpha=4/2=2\), для прочих \(\alpha\ge 2{,}5\). По принятому правилу берём слово, покрывающее больше новых гласных, то есть \emph{сыроежка}. Теперь \(C=\{\text{у},\text{я},\text{а},\text{е},\text{о},\text{ы}\}\), суммарная стоимость \(3+8=11\).

  \textbf{3-я итерация.} Наилучшие кандидаты — \emph{эвкалипт} с \(\alpha=8/2=4\) (новые \(\{\text{э},\text{и}\}\)) и \emph{липа} с \(\alpha=4/1=4\) (новая \(\{\text{и}\}\)). Выбираем \emph{эвкалипт}, так как он покрывает две новые гласные. Получаем \(C=\{\text{у},\text{я},\text{а},\text{е},\text{о},\text{ы},\text{э},\text{и}\}\), стоимость \(3+8+8=19\).

  \textbf{4-я итерация.} Остаются непокрытыми только \(\{\text{ё},\text{ю}\}\). Для \emph{опёнок} \(\alpha=6/1=6\), для \emph{берёза} \(\alpha=6/1=6\), для \emph{трюфель} \(\alpha=7/1=7\). По правилу разрешения равенств выбираем \emph{опёнок}. Теперь покрыта \(\text{ё}\), стоимость \(3+8+8+6=25\).

  \textbf{5-я итерация.} Добавляем \emph{трюфель}, чтобы покрыть \(\text{ю}\). Итого \(C=U\) и суммарная стоимость равна \(3+8+8+6+7=32\).




  \paragraph{Оптимальное решение и сравнение.}
  Буквы \(\text{ю}\) и \(\text{э}\) встречаются только в словах \emph{трюфель} и \emph{эвкалипт}, поэтому они обязательны в любом решении (стоимость \(7+8=15\)). Добавляя \emph{туя}  \(3\)), \emph{опёнок} (\(6\)) и \emph{рыжик} (\(5\)), получаем оптимальное решение \(\{\text{эвкалипт},\text{трюфель},\text{опёнок},\text{туя},\text{рыжик}\}\) суммарной стоимости \(29\).

  Таким образом, жадный алгоритм Хватала даёт решение стоимости \(32\), а оптимум равен \(29\).

  \section*{Задача 1.2. Достижение оценки алгоритма Хватала}

  \paragraph{Условие.}
  Построить пример задачи о покрытии множествами, на котором достигается оценка алгоритма Хватала \(H_n\), даже если вес каждого множества равен 1.

  \paragraph{Решение.}
  Оценка алгоритма Хватала утверждает, что стоимость найденного решения не превосходит \(H_n \cdot \text{OPT}\), где \(H_n = 1 + \frac{1}{2} + \ldots + \frac{1}{n}\) — гармоническое число. Для достижения этой оценки при единичных весах необходимо использовать различные веса множеств.

  \paragraph{Классическая конструкция.}
  Универсум \(U = \{e_1, e_2, \ldots, e_n\}\). Множества:
  \begin{itemize}
    \item \(S_i = \{e_i\}\) с весом \(c(S_i) = \frac{1}{n-i+1}\) для \(i = 1, 2, \ldots, n\);
    \item \(S_{\text{full}} = U\) с весом \(c(S_{\text{full}}) = 1 + \varepsilon\) для малого \(\varepsilon > 0\).
  \end{itemize}

  \paragraph{Работа алгоритма Хватала.}
  На итерации \(i\) при \(C = \{e_1, \ldots, e_{i-1}\}\) имеем:
  \begin{itemize}
    \item Для \(S_i\): \(\alpha(S_i) = \frac{1/(n-i+1)}{1} = \frac{1}{n-i+1}\);
    \item Для \(S_{\text{full}}\): \(\alpha(S_{\text{full}}) = \frac{1+\varepsilon}{n-i+1}\).
  \end{itemize}
  Минимум достигается у \(S_i\), поэтому алгоритм последовательно выбирает \(S_1, S_2, \ldots, S_n\).

  Стоимость жадного решения: \(\frac{1}{n} + \frac{1}{n-1} + \ldots + \frac{1}{1} = H_n\).

  Оптимальное решение: выбрать \(S_{\text{full}}\), стоимость \(1 + \varepsilon\).

  Отношение: \(\frac{H_n}{1+\varepsilon} \approx H_n \approx \ln n\).

  \paragraph{Попытка построения примера с единичными весами.}
  Рассмотрим пример с единичными весами множеств (но разными размерами).

  Универсум \(U = \{e_1, e_2, e_3, e_4\}\). Множества (все веса = 1):
  \begin{itemize}
    \item \(S_1 = \{e_1, e_2\}\);
    \item \(S_2 = \{e_2, e_4\}\);
    \item \(S_3 = \{e_1, e_3\}\).
  \end{itemize}

  Алгоритм Хватала выбирает все три множества \(S_1, S_2, S_3\) (стоимость 3), в то время как оптимум — это \(\{S_2, S_3\}\) (стоимость 2). хз как сделать продолжить идею...
  \section*{Задача 1.3. Уровневый алгоритм для ограниченной частоты}

  \paragraph{Условие.}
  Рассмотрим задачу о покрытии, где каждый элемент входит не более чем в 3 подмножества. Построить 3-приближённый уровневый алгоритм.

  \paragraph{Уровневый алгоритм.}
  \begin{enumerate}
    \item \(\text{Sol} \gets \varnothing\), \(i \gets 0\)
    \item \(w'(S_j) \gets c(S_j)\) для всех \(S_j \in S\) (инициализация остаточных весов)
    \item \(C_0 \gets \varnothing\) (множество покрытых элементов)
    \item \textbf{while} \(C_i \neq U\) \textbf{do}
    \item \quad \(\alpha \gets \min_{S_j} \dfrac{w'(S_j)}{|S_j \setminus C_i|}\) (минимальное отношение остаточного веса к числу непокрытых элементов)
    \item \quad \textbf{for each} \(S_j \in S\) \textbf{do}
    \item \quad\quad \(t_i(S_j) \gets \alpha \cdot |S_j \setminus C_i|\) (вычитаемая величина)
    \item \quad\quad \(w'(S_j) \gets w'(S_j) - t_i(S_j)\) (обновляем остаточный вес)
    \item \quad \textbf{end for}
    \item \quad \(W_i \gets \{S_j \mid w'(S_j) = 0\}\) (множества с нулевым остаточным весом)
    \item \quad \(\text{Sol} \gets \text{Sol} \cup W_i\) (добавляем в решение)
    \item \quad \(C_{i+1} \gets C_i \cup \bigcup_{S_j \in W_i} S_j\) (обновляем покрытые элементы)
    \item \quad \(i \gets i + 1\)
    \item \textbf{end while}
    \item \textbf{return} \(\text{Sol}\)
  \end{enumerate}

  \paragraph{Анализ алгоритма.}
  Определим для каждого элемента \(u \in U\) его вес \(y_u\) как сумму всех вычитаемых величин \(\alpha\) на итерациях, когда \(u\) был непокрыт. Тогда для любого множества \(S_j \in \text{Sol}\) выполнено \(\sum_{u \in S_j} y_u = c(S_j)\).

  \begin{lemma}
  Пусть каждый элемент входит не более чем в 3 множества. Тогда стоимость решения, найденного уровневым алгоритмом, не превосходит \(3 \cdot \text{OPT}\).
  \end{lemma}

  \paragraph{Доказательство.}
  Пусть \(\text{OPT}^*\) — оптимальное покрытие. Поскольку \(\text{OPT}^*\) покрывает все элементы, имеем:
  \[
  \sum_{S_j \in \text{OPT}^*} \sum_{u \in S_j} y_u \geq \sum_{u \in U} y_u.
  \]
  Следовательно, \(\text{OPT} = \sum_{S_j \in \text{OPT}^*} c(S_j) \geq \sum_{u \in U} y_u\).

  Стоимость найденного решения:
  \[
  \text{ALG} = \sum_{S_j \in \text{Sol}} c(S_j) = \sum_{S_j \in \text{Sol}} \sum_{u \in S_j} y_u = \sum_{u \in U} \sum_{S_j \in \text{Sol}: u \in S_j} y_u.
  \]
  Поскольку каждый элемент \(u\) входит не более чем в 3 множества из \(S\), он входит не более чем в 3 множества из \(\text{Sol}\). Следовательно:
  \[
  \sum_{S_j \in \text{Sol}: u \in S_j} y_u \leq 3 \cdot y_u.
  \]
  Отсюда:
  \[
  \text{ALG} \leq \sum_{u \in U} 3 \cdot y_u = 3 \sum_{u \in U} y_u \leq 3 \cdot \text{OPT}.
  \]
  Таким образом, алгоритм является 3-приближённым.

  \section*{Задача 1.4. Обратное вершинное покрытие в полном антисимметрическом графе}

  \paragraph{Условие.}
  Полным антисимметрическим графом называется ориентированный граф \(G = (V, E)\), в котором для любой пары вершин \(u, v \in V\) существует ровно одна из дуг \((u, v) \in E\) либо \((v, u) \in E\). Обратным вершинным покрытием (Feedback Vertex Set, FVS) в \(G\) называется множество вершин, после удаления которых граф становится ациклическим. Предложить 3-приближённый алгоритм для поиска обратного вершинного покрытия.

  \begin{lemma}
  Если полный антисимметрический граф не содержит циклов длины 3, то он ациклический.
  \end{lemma}

  \paragraph{Доказательство.}
  Предположим противное: граф не содержит треугольников (циклов длины 3), но содержит цикл длины \(k \geq 4\):
  \[
  v_1 \to v_2 \to v_3 \to \ldots \to v_k \to v_1.
  \]
  Рассмотрим вершины \(v_1, v_2, v_3\). Из цикла имеем дуги \(v_1 \to v_2\) и \(v_2 \to v_3\). Поскольку граф полный, между \(v_1\) и \(v_3\) существует дуга.

  Если \(v_1 \to v_3\), то вместе с дугой \(v_3 \to v_4 \to \ldots \to v_k \to v_1\) получаем более короткий цикл. Продолжая сокращать, придём к треугольнику.

  Если \(v_3 \to v_1\), то \(v_1 \to v_2 \to v_3 \to v_1\) образует треугольник — противоречие.

  Следовательно, в полном антисимметрическом графе без треугольников нет циклов.

  \paragraph{Сведение к покрытию треугольников.}
  Из леммы следует, что для получения ациклического графа достаточно удалить вершины так, чтобы разбить все треугольники. Задача сводится к задаче о покрытии множествами:
  \begin{itemize}
    \item Универсум \(U\) — множество всех треугольников в графе.
    \item Для каждой вершины \(v \in V\) определим множество \(S_v\) — все треугольники, содержащие \(v\).
    \item Требуется найти минимальное множество вершин, покрывающее все треугольники.
  \end{itemize}

  \paragraph{Ключевое наблюдение.}
  Каждый треугольник содержит ровно 3 вершины, следовательно, каждый элемент универсума \(U\) (треугольник) входит ровно в 3 множества \(S_{v_1}, S_{v_2}, S_{v_3}\), где \(v_1, v_2, v_3\) — вершины треугольника. Это задача о покрытии с частотой 3.

  \paragraph{Алгоритм.}
  \begin{enumerate}
    \item Найти все треугольники в графе (перебор всех троек вершин за \(O(n^3)\)).
    \item Применить уровневый алгоритм из задачи 1.3 для покрытия треугольников.
  \end{enumerate}

  По задаче 1.3 уровневый алгоритм для задачи о покрытии с частотой 3 является 3-приближённым. Следовательно, найденное обратное вершинное покрытие имеет размер не более \(3 \cdot \text{OPT}\).

\end{document}


